\section{Conclusion}

Completion of this laboratory work has provided practical understanding of analog signal
acquisition and conditioning techniques for embedded sensor systems, and reinforced the
use of FreeRTOS synchronization primitives for safe inter-task communication.

This sixth laboratory work covered the complete sensor pipeline from raw data acquisition
through multi-stage signal conditioning to structured reporting. The key takeaways are:

\begin{enumerate}
  \item Implementing \textbf{periodic sensor acquisition} with the DHT11 in FreeRTOS
    context -- the sensor's minimum 2\,second interval between reads defines the
    acquisition period, and the active sensor is selected from a configurable catalogue,
    allowing the system to be extended with additional sensor types without modifying
    the task logic.
  \item Applying a \textbf{three-stage signal conditioning pipeline}:
    \begin{itemize}
      \item \textbf{Saturation} clamps the raw value to the sensor's physical range
        ($0$--$50\,^\circ$C), rejecting impossible readings at the input boundary.
      \item \textbf{Median filter} (window = 5) removes salt-and-pepper noise --
        isolated outlier spikes are eliminated without distorting the underlying
        temperature trend, as the median is inherently robust to up to
        $\lfloor n/2 \rfloor$ outliers in the window.
      \item \textbf{Exponential Moving Average} ($\alpha = 0.2$) provides weighted
        smoothing where recent samples contribute 20\% and the accumulated history
        contributes 80\%, yielding a smooth output suitable for display while
        maintaining adequate responsiveness to real temperature changes.
    \end{itemize}
  \item Evaluating \textbf{alert thresholds with hysteresis} (ON $\geq 28\,^\circ$C,
    OFF $\leq 24\,^\circ$C) on the median-filtered value rather than the heavily
    smoothed EMA, ensuring that the alert responds promptly to actual temperature
    crossings while the dead band prevents chattering near the threshold.
  \item Structuring the firmware into three distinct \textbf{FreeRTOS tasks} with
    well-defined roles and priorities: Acquisition (priority 3) runs most frequently and
    ensures fresh data, Conditioning (priority 2) is event-driven via semaphore and
    reacts immediately to new samples, and Reporter (priority 1) provides periodic
    human-readable output without interfering with the control loop.
  \item Using a \textbf{mutex-protected shared data store} (\texttt{SensorData}) as
    the single source of truth for all sensor state -- raw, filtered, and averaged
    values along with validity and alert state -- preventing data races when multiple
    tasks read or write concurrently. All configurable pipeline parameters (acquisition
    period, saturation limits, median window, EMA coefficient, thresholds) are
    centralized in the \texttt{SensorData.h} header.
  \item Organizing the project into \textbf{three hardware-abstraction layers} (MCAL,
    ECAL, SRV) with PlatformIO, keeping all hardware-specific code in the MCAL layer,
    declaring all pin assignments in ECAL driver headers for easy circuit reference,
    and making the service layer fully portable across different MCU platforms.
\end{enumerate}

This laboratory work demonstrates how a pipeline of classical signal processing
techniques -- saturation, median filtering, and exponential moving averaging --
combined with FreeRTOS synchronization primitives, transforms a noisy raw sensor
stream into a clean, reliable signal suitable for both display and alert generation
in an embedded system.
